\documentclass{beamer}

% ---------- Thème et couleurs ----------
\usetheme{Madrid}
\usecolortheme{beaver}
\setbeamertemplate{navigation symbols}{}
\setbeamertemplate{footline}[frame number]

% ---------- Packages ----------
\usepackage[utf8]{inputenc}
\usepackage[T1]{fontenc}
\usepackage[french]{babel}
\usepackage{graphicx}
\usepackage{array}
\usepackage{booktabs}
\usepackage{xcolor}

% ---------- Titre ----------
\title{\textbf{Workflows IoT Intelligents}}
\subtitle{Privacy-Aware Distributed CNN \& MARL-Based IoT Security System}
\author{Présenté par : SOULAIMANE Ait Ahmed Oulhaj}
\institute{Master AIDC — FST Béni Mellal}
\date{\today}

% ==========================================================
\begin{document}

% ---------- Page de titre ----------
\begin{frame}
  \titlepage
\end{frame}

% ---------- Introduction ----------
\begin{frame}{Introduction}
\begin{itemize}
    \item L’\textbf{Internet des Objets (IoT)} et l’\textbf’intelligence artificielle (IA) ont profondément transformé notre manière de collecter, traiter et analyser les données.
    \item Chaque jour, des milliards de capteurs, caméras et dispositifs intelligents produisent d’énormes volumes de données hétérogènes et sensibles.
    \item Ces données sont souvent exploitées par des modèles profonds tels que les \textbf{réseaux de neurones convolutionnels (CNN)} pour des tâches d’analyse avancées :
    \begin{itemize}
        \item vision par ordinateur (surveillance, reconnaissance faciale),
        \item analyse médicale (IRM, radiologie),
        \item systèmes intelligents (véhicules autonomes, smart cities).
    \end{itemize}
\end{itemize}
\end{frame}

\begin{frame}{Problématique}
\begin{itemize}
    \item Traditionnellement, les calculs d’inférence sont délégués à des \textbf{serveurs cloud distants} capables de gérer la lourde charge computationnelle.
    \item Cependant, cette centralisation pose deux limites majeures :
    \begin{enumerate}
        \item \textbf{Performance :} la latence due à la transmission des données vers le cloud compromet les applications temps réel.
        \item \textbf{Confidentialité :} les données brutes (images, signaux, logs) quittent leur source et exposent des informations sensibles.
    \end{enumerate}
    \item Face à ces défis, l’\textbf{Edge Computing} et l’\textbf{apprentissage distribué} émergent comme solutions pour rapprocher l’IA des sources de données.
\end{itemize}
\end{frame}

\begin{frame}{Motivation et idée clé}
\begin{itemize}
    \item Pour répondre simultanément aux besoins de \textbf{confidentialité}, de \textbf{scalabilité} et de \textbf{rapidité},
    les chercheurs ont proposé une \textbf{architecture distribuée et coopérative} combinant deux paradigmes :
    \begin{itemize}
        \item  \textbf{Privacy-Aware Distributed CNN} : 
        distribution des couches de CNN entre plusieurs appareils IoT pour effectuer le calcul près de la source tout en brouillant les données sensibles.
        \item  \textbf{MARL-Based IoT Security System} : 
        utilisation de l’apprentissage par renforcement multi-agent (MARL) pour orchestrer la coopération entre nœuds et protéger le réseau contre les attaques.
    \end{itemize}
    \item Cette combinaison vise à créer un système \textbf{intelligent, décentralisé et respectueux de la vie privée},
    capable de traiter des tâches complexes sans dépendre d’un serveur central.
\end{itemize}
\end{frame}

\begin{frame}{Objectifs du travail}
\begin{itemize}
    \item Développer une \textbf{stratégie de distribution et de coopération} pour exécuter des tâches d’inférence CNN en périphérie du réseau.
    \item Garantir :
    \begin{itemize}
        \item  la \textbf{confidentialité} des données locales ;
        \item  la \textbf{réduction de la latence} d’exécution ;
        \item  la \textbf{collaboration efficace} entre entités IoT hétérogènes.
    \end{itemize}
    \item Exploiter le \textbf{MARL compétitif et coopératif} pour permettre à chaque agent IoT :
    \begin{itemize}
        \item d’apprendre à répartir ses tâches CNN de manière autonome,
        \item tout en tenant compte des contraintes globales du réseau (ressources, bande passante, confidentialité).
    \end{itemize}
\end{itemize}
\end{frame}

\begin{frame}{Contexte général}
\begin{itemize}
    \item Dans cet environnement, chaque entité (ex. hôpital, aéroport, banque) agit comme un \textbf{agent intelligent} :
    \begin{itemize}
        \item Elle conserve ses données privées ;
        \item Elle collabore avec d’autres nœuds IoT pour effectuer une partie du calcul CNN ;
        \item Elle optimise ses choix à l’aide d’un apprentissage par renforcement.
    \end{itemize}
    \item Ce travail établit ainsi un \textbf{lien entre apprentissage profond et apprentissage multi-agent}, pour concevoir une IA :
    \begin{itemize}
        \item plus \textbf{rapide} (réduction de latence),
        \item plus \textbf{sécurisée} (protection des données),
        \item et plus \textbf{coopérative} (gestion partagée des ressources).
    \end{itemize}
\end{itemize}
\end{frame}



% ==========================================================
% ---------- Privacy-Aware Distributed CNN ----------
\begin{frame}{Workflow : Privacy-Aware Distributed CNN}
\textbf{Objectif :} Réaliser des tâches d’analyse (détection d’objets, IRM, reconnaissance faciale) sans divulguer les données originales.
\begin{enumerate}
    \item \textbf{Collecte locale :} chaque appareil IoT capture les données brutes (images, signaux...).
    \item \textbf{Pré-traitement local :} premières couches CNN appliquées localement → création de \textit{feature maps floutées}.
    \item \textbf{Distribution du calcul :} les features sont réparties entre plusieurs nœuds IoT selon leurs ressources.
    \item \textbf{Calcul collaboratif :} chaque nœud exécute quelques couches CNN.
    \item \textbf{Agrégation et sortie locale :} les résultats partiels sont réunis, la dernière couche s’exécute localement.
\end{enumerate}
\end{frame}

% ---------- Avantages Privacy-Aware CNN ----------
\begin{frame}{Avantages — Privacy-Aware Distributed CNN}
\begin{itemize}
    \item ️ Protection forte de la vie privée (aucune donnée brute transmise)
    \item ⚙️ Utilisation optimale des ressources IoT
    \item  Scalabilité grâce à l’exécution distribuée
    \item  Résilience aux attaques : données floutées
    \item  Réduction de la charge du cloud
\end{itemize}
\end{frame}

% ---------- Inconvénients Privacy-Aware CNN ----------
\begin{frame}{Inconvénients — Privacy-Aware Distributed CNN}
\begin{itemize}
    \item  Synchronisation complexe entre nœuds
    \item  Dépendance à la connectivité réseau
    \item  Consommation énergétique accrue
    \item  Découpage du CNN difficile à optimiser
    \item  Modèle statique (pas d’apprentissage continu)
\end{itemize}
\end{frame}

% ==========================================================
% ---------- MARL-Based IoT Security System ----------
\begin{frame}{Workflow : MARL-Based IoT Security System}
\textbf{Objectif :} Détection et réponse collective aux attaques dans un réseau IoT.
\begin{enumerate}
    \item \textbf{Observation locale :} chaque agent surveille son environnement (trafic, logs...).
    \item \textbf{Détection et action :} identification d’anomalies et actions locales (isolation, alerte...).
    \item \textbf{Apprentissage local (RL) :} mise à jour de la politique $\pi_i(a_i|s_i)$.
    \item \textbf{Agrégation fédérée (FedAvg) :} les poids $\theta_i$ sont envoyés au coordinateur.
    \item \textbf{CTDE :} le serveur critique évalue et agrège les politiques globales.
    \item \textbf{Exécution décentralisée :} chaque agent agit de manière autonome.
\end{enumerate}
\end{frame}

% ---------- Avantages MARL ----------
\begin{frame}{Avantages — MARL-Based IoT Security}
\begin{itemize}
    \item  Apprentissage collectif et adaptatif
    \item ️ Défense autonome contre les attaques
    \item  Confidentialité préservée (échange de poids, pas de données)
    \item ⚙️ Résilience du réseau : pas de point de défaillance unique
    \item  Amélioration continue grâce au RL fédéré
\end{itemize}
\end{frame}

% ---------- Inconvénients MARL ----------
\begin{frame}{Inconvénients — MARL-Based IoT Security}
\begin{itemize}
    \item  Coordination complexe entre agents
    \item  Communication fréquente (poids/politiques)
    \item ⚖️ Risque de modèles biaisés ou corrompus
    \item  Coût computationnel élevé (RL)
    \item  Vulnérabilité au \textit{poisoning} fédéré
\end{itemize}
\end{frame}

% ==========================================================
% ---------- Comparaison ----------
\begin{frame}{Comparaison des deux approches}
\small
\begin{tabular}{>{\bfseries}m{3.2cm} m{4cm} m{4cm}}
\toprule
\textbf{Élément} & \textbf{Privacy-Aware CNN} & \textbf{MARL-Based IoT Security} \\
\midrule
Objectif & Préserver la vie privée des données & Détecter et prévenir les attaques \\
Type d’apprentissage & Inférence distribuée & RL fédéré multi-agent \\
Données partagées & Features floutées & Poids du modèle \\
Contrôle central & Source CNN locale & Coordinateur fédéré (CTDE) \\
Résultat final & Sortie locale confidentielle & Modèle global robuste \\
Type de sécurité & Confidentialité & Cybersécurité \\
Adaptativité & Non & Oui (auto-apprentissage) \\
\bottomrule
\end{tabular}
\end{frame}

% ==========================================================

 \begin{frame}{Idée développée après l’introduction}
\begin{itemize}
    \item Après l’introduction, l’article développe trois axes majeurs visant à concilier \textbf{performance}, \textbf{confidentialité} et \textbf{intelligence distribuée} dans les réseaux IoT.
\end{itemize}

\vspace{0.3cm}
\begin{block}{ (A) Distribution de CNN sans serveur central}
\begin{itemize}
    \item Chaque entité (ex. hôpital, aéroport, banque) conserve ses données localement.
    \item Les couches du CNN sont \textbf{découpées et distribuées} entre plusieurs dispositifs IoT voisins.
    \item Objectif : \textbf{réduire la latence d’inférence} tout en garantissant qu’aucune donnée brute ne quitte la source.
\end{itemize}
\end{block}
\end{frame}

\begin{frame}{Idée développée après l’introduction}

\begin{block}{ (B) Formulation en jeu multi-agent}
\begin{itemize}
    \item Chaque entité est modélisée comme un \textbf{agent autonome} cherchant à optimiser ses propres objectifs :
    \begin{itemize}
        \item rapidité du traitement,
        \item préservation de la vie privée.
    \end{itemize}
    \item Les agents interagissent dans un \textbf{environnement partagé} et compétitif, formant un \textbf{jeu stochastique}.
\end{itemize}
\end{block}
\end{frame}

\begin{frame}{Idée développée après l’introduction}
\begin{block}{ (C) Utilisation du MARL (Competitive + Cooperative)}
\begin{itemize}
    \item Chaque agent apprend dynamiquement à \textbf{distribuer ses tâches CNN} selon :
    \begin{itemize}
        \item ses besoins de confidentialité ;
        \item l’état global des ressources disponibles (CPU, mémoire, bande passante) ;
        \item les actions des autres agents du réseau.
    \end{itemize}
    \item Le \textbf{MARL (Multi-Agent Reinforcement Learning)} permet au système de :
    \begin{itemize}
        \item s’adapter en temps réel sans supervision centrale ;
        \item équilibrer la compétition et la coopération ;
        \item atteindre un fonctionnement globalement optimal et sécurisé.
    \end{itemize}
\end{itemize}
\end{block}
\end{frame}

% ---------- Conclusion ----------
\begin{frame}{Conclusion}
\begin{itemize}
    \item  Le \textbf{Privacy-Aware Distributed CNN} protège les données pendant le calcul.
    \item  Le \textbf{MARL-Based IoT Security System} protège le réseau contre les intrusions.
    \item  Les deux approches peuvent être \textbf{complémentaires} :
    \begin{itemize}
        \item CNN distribué pour le traitement privé.
        \item MARL fédéré pour la sécurité adaptative.
    \end{itemize}
    \item  Vers un futur \textbf{IoT intelligent, sécurisé et respectueux de la vie privée}.
\end{itemize}
\end{frame}

% ==========================================================
\begin{frame}{Fonctionnement du Privacy-Aware Distributed CNN}
\textbf{Principe général :}
\begin{itemize}
    \item Exécuter un \textbf{réseau CNN} de manière \textbf{distribuée} sur plusieurs appareils IoT.
    \item Préserver la confidentialité des données en évitant tout transfert d’images ou de signaux bruts.
    \item Chaque dispositif IoT contribue au calcul selon ses capacités (CPU, mémoire, bande passante).
\end{itemize}

\vspace{0.4cm}
\textbf{Objectif :} réaliser des tâches d’analyse (vision, IRM, détection d’objets) en garantissant :
\begin{itemize}
    \item  \textbf{Vie privée} : les données originales ne quittent jamais la source.
    \item  \textbf{Performance} : calcul proche de la source pour réduire la latence.
    \item  \textbf{Collaboration} : plusieurs nœuds IoT coopèrent au traitement.
\end{itemize}
\end{frame}

% ==========================================================
\begin{frame}{Étapes détaillées du fonctionnement}
\begin{enumerate}
    \item \textbf{Collecte locale des données :}  
    Chaque appareil IoT capture des données (image, signal, vidéo) et les conserve localement.

    \item \textbf{Pré-traitement local (brouillage) :}  
    Les premières couches du CNN (convolution + pooling) sont exécutées localement pour générer des \textbf{feature maps floutées}, ne contenant plus d’informations identifiables.

    \item \textbf{Fragmentation et distribution :}  
    Les feature maps sont \textbf{fragmentées} et distribuées entre plusieurs nœuds IoT voisins selon leurs ressources disponibles :
    \[
    R_i = f(\bar{m_i}, \bar{c_i}, \bar{b_i})
    \]
    où $\bar{m_i}$, $\bar{c_i}$ et $\bar{b_i}$ désignent respectivement la mémoire, la puissance CPU et la bande passante du nœud $i$.

    \item \textbf{Calcul collaboratif :}  
    Chaque nœud exécute une partie du CNN (opérations intermédiaires) et renvoie ses résultats partiels.
\end{enumerate}
\end{frame}

% ==========================================================
\begin{frame}{Fin du processus et sécurité}
\begin{enumerate}
    \setcounter{enumi}{4}
    \item \textbf{Sécurité et anonymisation :}  
    Les échanges sont chiffrés (TLS/SSL) et les métadonnées sensibles supprimées.  
    Même en cas de compromission, un nœud ne possède qu’une portion floutée des données.

    \item \textbf{Agrégation des résultats :}  
    Le dispositif source rassemble les résultats partiels provenant des nœuds pour reconstituer le flux d’inférence complet.

    \item \textbf{Dernières couches CNN :}  
    Les dernières couches (Fully Connected + Softmax) s’exécutent localement pour produire la \textbf{sortie finale confidentielle}.
\end{enumerate}

\vspace{0.4cm}
\textbf{Schéma du flux :}
\begin{center}
\scriptsize
\texttt{[Donnée locale]} → \texttt{[CNN local (floutage)]} → \texttt{[Feature maps distribuées]} → \texttt{[Calcul IoT collaboratif]} → \texttt{[Agrégation]} → \texttt{[Résultat final local]}
\end{center}
\end{frame}

% ==========================================================
\begin{frame}{Avantages du Privacy-Aware Distributed CNN}
\begin{itemize}
    \item ️ \textbf{Protection de la vie privée} : aucune donnée brute n’est transmise.
    \item ⚙️ \textbf{Utilisation optimale des ressources IoT} : chaque nœud travaille selon sa capacité.
    \item  \textbf{Architecture scalable} : plus de nœuds = plus de puissance collective.
    \item  \textbf{Résilience aux attaques} : un nœud compromis ne peut pas reconstruire les données originales.
    \item  \textbf{Faible latence} : traitement en périphérie, idéal pour le temps réel.
\end{itemize}

\vspace{0.4cm}
\begin{block}{En résumé}
Le Privacy-Aware Distributed CNN repose sur trois piliers :
\begin{enumerate}
    \item Exécution locale partielle (brouillage des données),
    \item Distribution intelligente entre nœuds,
    \item Agrégation sécurisée et résultat final local.
\end{enumerate}
\end{block}
\end{frame}

% ==========================================================
\begin{frame}{Principe du Privacy-Aware Distributed CNN}
\begin{itemize}
    \item L’idée clé est de distribuer le calcul d’un \textbf{réseau CNN} entre plusieurs appareils IoT tout en \textbf{préservant la confidentialité} des données.
    \item Les \textbf{premières couches CNN} (convolution, pooling) sont exécutées localement sur l’appareil source pour produire des \textbf{feature maps floutées}.
    \item Ces feature maps sont ensuite \textbf{fragmentées} et envoyées à différents nœuds IoT pour traitement parallèle.
    \item Enfin, les résultats partiels reviennent vers le dispositif d’origine pour exécuter les \textbf{dernières couches CNN} (Fully Connected + Softmax) et obtenir la classification finale.
\end{itemize}

\vspace{0.3cm}
\textbf{But :} Réduire la latence et la charge du cloud, sans jamais exposer les données brutes.
\end{frame}

% ==========================================================
\begin{frame}{Distribution des Feature Maps entre Nœuds IoT}
\begin{itemize}
    \item Après la première convolution, une image produit plusieurs \textbf{feature maps} :
    \[
    F = \{F_1, F_2, F_3, \dots, F_n\}
    \]
    \item Chaque feature map représente une caractéristique visuelle (bords, textures, formes…).
    \item Ces feature maps sont réparties entre différents nœuds :
\end{itemize}

\vspace{0.2cm}
\begin{center}
\scriptsize
\begin{tabular}{c c c c}
\textbf{Nœud IoT 1} & \textbf{Nœud IoT 2} & \textbf{Nœud IoT 3} & \textbf{Nœud IoT 4} \\
\hline
$F_1$ & $F_2$ & $F_3$ & $F_4$ \\
$\downarrow$ & $\downarrow$ & $\downarrow$ & $\downarrow$ \\
Conv + Pool & Conv + Pool & Conv + Pool & Conv + Pool \\
$\downarrow$ & $\downarrow$ & $\downarrow$ & $\downarrow$ \\
$\hat{F}_1$ & $\hat{F}_2$ & $\hat{F}_3$ & $\hat{F}_4$ \\
\end{tabular}
\end{center}

\vspace{0.3cm}
\begin{itemize}
    \item Les résultats partiels $\hat{F}_i$ sont renvoyés et \textbf{agrégés localement} :
    \[
    \text{Concat}(\hat{F}_1, \hat{F}_2, \hat{F}_3, \hat{F}_4) \rightarrow \text{Sortie finale}
    \]
\end{itemize}
\end{frame}

% ==========================================================
\begin{frame}{Coordination entre Nœuds et Rôle du MARL}
\textbf{Deux scénarios possibles :}
\begin{block}{1️⃣ Système homogène}
\begin{itemize}
    \item Tous les nœuds IoT participent au même modèle CNN et à la même tâche (ex : détection d’objets).
    \item Chacun traite une portion différente des feature maps, mais les poids du modèle sont identiques.
\end{itemize}
\end{block}

\begin{block}{2️⃣ Système hétérogène}
\begin{itemize}
    \item Plusieurs entités (ex. hôpital, aéroport, banque) ont leurs propres modèles CNN.
    \item Elles partagent les mêmes ressources IoT, gérées via un système \textbf{MARL (Multi-Agent Reinforcement Learning)}.
    \item Le MARL apprend à :
    \begin{itemize}
        \item allouer les ressources selon les besoins de chaque agent ;
        \item équilibrer la charge pour minimiser la latence ;
        \item préserver la confidentialité globale.
    \end{itemize}
\end{itemize}
\end{block}
\end{frame}

% ==========================================================
\begin{frame}{Synthèse et Sécurité du Processus}
\begin{itemize}
    \item Même si un nœud est compromis :
    \begin{itemize}
        \item il ne détient qu’une \textbf{portion floutée} des feature maps ;
        \item il ne peut pas reconstruire l’entrée originale ;
        \item il ne connaît ni les autres portions ni la structure complète du CNN.
    \end{itemize}
    \item Les communications sont chiffrées (TLS/SSL) et les métadonnées sensibles supprimées.
    \item Le dispositif source réalise la \textbf{classification finale localement}, garantissant la confidentialité du résultat.
\end{itemize}

\vspace{0.3cm}
\begin{center}
\scriptsize
\texttt{[Image locale]} → \texttt{[Conv + Pool (local)]} → \texttt{[Feature maps distribuées]} → \texttt{[Calcul IoT parallèle]} → \texttt{[Agrégation + Sortie locale]}
\end{center}

\vspace{0.2cm}
\begin{block}{En résumé}
Le Privacy-Aware Distributed CNN combine :
\begin{itemize}
    \item  \textbf{Vie privée} — données jamais exposées.
    \item ⚙️ \textbf{Calcul distribué} — répartition adaptative selon les ressources.
    \item  \textbf{Apprentissage coopératif} — coordination via MARL.
\end{itemize}
\end{block}
\end{frame}

% ==========================================================
\begin{frame}{Principe du MARL-Based IoT Security System}
\begin{itemize}
    \item Le \textbf{Multi-Agent Reinforcement Learning (MARL)} permet à plusieurs appareils IoT d’apprendre collectivement à \textbf{détecter et contrer des attaques}.
    \item Chaque appareil (caméra, capteur, routeur…) agit comme un \textbf{agent intelligent} :
    \begin{itemize}
        \item il observe son environnement local ;
        \item il prend des décisions autonomes (actions) ;
        \item il reçoit des récompenses selon l’efficacité de ses réponses ;
        \item il partage uniquement ses paramètres de modèle, pas ses données.
    \end{itemize}
    \item Objectif global : créer une \textbf{défense collective, distribuée et auto-adaptative} du réseau IoT.
\end{itemize}
\end{frame}

% ==========================================================
\begin{frame}{Étapes 1-2 : Observation et Action Locale}
\textbf{1️⃣ Observation locale :}
\begin{itemize}
    \item Chaque agent surveille son propre environnement :
    \begin{itemize}
        \item trafic réseau, logs, métriques CPU ;
        \item activité des capteurs ou caméras ;
        \item anomalies dans le flux de données.
    \end{itemize}
    \item Cet état local est noté $s_i(t)$.
\end{itemize}

\vspace{0.3cm}
\textbf{2️⃣ Détection et action locale :}
\begin{itemize}
    \item L’agent choisit une action $a_i = \pi_i(s_i)$ selon sa politique apprise :
    \begin{itemize}
        \item isoler un nœud compromis ;
        \item bloquer une adresse IP ;
        \item alerter les voisins.
    \end{itemize}
    \item Une récompense $r_i$ est attribuée selon la pertinence de l’action :
    \[
    r_i > 0 \Rightarrow \text{attaque évitée}, \quad r_i < 0 \Rightarrow \text{alerte inutile.}
    \]
\end{itemize}
\end{frame}

% ==========================================================
\begin{frame}{Étape 3 : Entraînement local et Partage fédéré}
\textbf{3️⃣ Apprentissage local (RL individuel) :}
\begin{itemize}
    \item Chaque agent met à jour sa fonction de valeur $Q_i(s_i,a_i)$ via Q-learning :
    \[
    Q_i(s_i,a_i) \leftarrow Q_i(s_i,a_i) + \alpha [r_i + \gamma \max_{a'} Q_i(s'_i,a') - Q_i(s_i,a_i)]
    \]
    \item L’agent améliore ainsi sa capacité à réagir face aux attaques futures.
\end{itemize}

\vspace{0.3cm}
\textbf{4️⃣ Partage fédéré (Federated MARL) :}
\begin{itemize}
    \item Les agents n’envoient pas leurs données, mais seulement leurs \textbf{poids de modèle} $\theta_i$.
    \item Le coordinateur fédéré effectue une agrégation moyenne :
    \[
    \theta_{global} = \frac{1}{N} \sum_{i=1}^{N} \theta_i
    \]
    \item Le modèle global devient ainsi plus robuste et général, sans compromettre la confidentialité.
\end{itemize}
\end{frame}

% ==========================================================
\begin{frame}{Étapes 5-6 : CTDE et Exécution Décentralisée}
\textbf{5️⃣ Entraînement centralisé (CTDE) :}
\begin{itemize}
    \item Le coordinateur agit comme un \textbf{critique global} :
    \begin{itemize}
        \item il évalue la performance collective des politiques ;
        \item il pondère les contributions selon un niveau de confiance $w_i$ (Trust Layer) ;
        \item il met à jour un modèle global partagé.
    \end{itemize}
    \item Nouveaux paramètres :
    \[
    \theta_{new} = \sum_i w_i \cdot \theta_i
    \]
\end{itemize}

\vspace{0.3cm}
\textbf{6️⃣ Exécution décentralisée :}
\begin{itemize}
    \item Chaque agent reçoit la politique globale $\pi_{global}$.
    \item Il agit ensuite de façon autonome en temps réel :
    \begin{itemize}
        \item observation → décision → action → récompense ;
        \item apprentissage continu et mise à jour du modèle.
    \end{itemize}
\end{itemize}
\end{frame}

% ==========================================================
\begin{frame}{Synthèse du MARL-Based IoT Security System}
\begin{center}
\scriptsize
\texttt{
[Agent IoT local]
  ↓\\
Observation (état s\_i)
  ↓\\
Détection d’anomalie (CNN/LSTM/RL)
  ↓\\
Action (a\_i) → Récompense (r\_i)
  ↓\\
Mise à jour du modèle local
  ↓\\
↑  Partage des poids (θ\_i)
│\\
└──→ Coordinateur fédéré
        ↓\\
 Agrégation (CTDE → modèle global)
        ↓\\
 Redistribution aux agents
        ↓\\
[Exécution décentralisée en temps réel]
}
\end{center}

\vspace{0.3cm}
\begin{block}{Avantages}
\begin{itemize}
    \item  \textbf{Apprentissage adaptatif} — le système s’améliore à chaque itération.
    \item  \textbf{Préservation de la vie privée} — aucun partage de données brutes.
    \item  \textbf{Résilience et coopération} — même si un nœud échoue, les autres maintiennent la défense.
    \item  \textbf{Réactivité locale} — détection et réaction instantanées.
\end{itemize}
\end{block}
\end{frame}

\end{document}